\documentclass[a4paper,12pt]{article}

\usepackage[T2A]{fontenc}
\usepackage[utf8]{inputenc}
\usepackage[russian]{babel}
\usepackage{amsmath,amssymb,amsfonts}
\usepackage{graphicx}
\usepackage{float}
\usepackage{caption}
\usepackage{subcaption}
\usepackage{geometry}
\usepackage{setspace}
\usepackage{indentfirst}
\usepackage{enumitem}
\usepackage{booktabs}
\usepackage{algorithm}
\usepackage{algpseudocode}
\usepackage{hyperref}
\usepackage{amsthm}
\usepackage{thmtools}
\usepackage{xcolor}

\geometry{left=3cm, right=1.5cm, top=2cm, bottom=2cm}
\onehalfspacing
\setlength{\parindent}{1.25cm}
\setlength{\parskip}{0pt}
\hypersetup{
    colorlinks=true,
    linkcolor=blue,
    filecolor=magenta,
    citecolor=red,
    urlcolor=cyan,
    pdftitle={Курсовой проект},
    pdfpagemode=FullScreen,
}

\declaretheoremstyle[
    spaceabove=6pt, spacebelow=6pt,
    headfont=\normalfont\bfseries,
    notefont=\mdseries, notebraces={(}{)},
    bodyfont=\normalfont,
    postheadspace=1em,
    headpunct={}
]{mystyle}

\declaretheorem[style=mystyle, name=Определение, numberwithin=section]{definition}
\declaretheorem[style=mystyle, name=Теорема, sibling=definition]{theorem}
\declaretheorem[style=mystyle, name=Утверждение, sibling=definition]{proposition}
\declaretheorem[style=mystyle, name=Замечание, numbered=no]{remark}
\renewcommand{\theHdefinition}{\thesection.\arabic{definition}}

\floatname{algorithm}{Алгоритм}
\algrenewcommand\algorithmicrequire{\textbf{Вход:}}
\algrenewcommand\algorithmicensure{\textbf{Выход:}}

\newcommand{\R}{\mathbb{R}}
\newcommand{\norm}[1]{\left\lVert #1 \right\rVert}
\newcommand{\fro}[1]{\left\lVert #1 \right\rVert_{\mathrm F}}
\newcommand{\nuc}[1]{\left\lVert #1 \right\rVert_*}
\newcommand{\tr}{\operatorname{tr}}
\newcommand{\rank}{\operatorname{rank}}
\newcommand{\grad}{\operatorname{grad}}
\renewcommand{\Proj}{\mathcal{P}}
\newcommand{\Mcal}{\mathcal{M}}
\newcommand{\Retr}{\operatorname{Retr}}
\newcommand{\argmin}{\operatorname*{arg\,min}}

\begin{document}

\begin{titlepage}
    \centering
    {\large\bfseries Национальный исследовательский университет\\
    «Высшая школа экономики»\par}
    \vspace{1cm}
    {\large Факультет экономических наук, Факультет компьютерных наук\par}
    {\large Образовательные программы «Экономика и анализ данных», «Прикладная математика и информатика»\par}
    \vspace{2cm}
    {\Large\bfseries КУРСОВОЙ ПРОЕКТ\par}
    \vspace{1cm}
    {\Large\bfseries Геометрические методы на многообразии матриц фиксированного ранга для заполнения матриц\par}
    \vspace{2cm}
    \begin{flushright}
        \large
        \textbf{Выполнили:}\\
        Студент 2 курса, группы БЭАД245\\
        Аляутдинов Карим Маратович,\\
        Студент 2 курса, группы БПМИИ245\\
        Исмаилов Магомед Айдынович.\\
        \vspace{0.5cm}
        \textbf{Научный руководитель:}\\
        профессор, д.ф.-м.н.\\
        Широков Дмитрий Сергеевич\\
    \end{flushright}
    \vfill
    {\large Москва, 2026\par}
\end{titlepage}

\tableofcontents
\newpage

\section*{Введение}
\phantomsection
\addcontentsline{toc}{section}{Введение}

\textbf{Актуальность темы.} Задача восстановления пропусков в матрицах возникает во многих прикладных задачах: в рекомендательных системах, обработке изображений, восстановлении панельных экономических и социальных данных. Часто именно восстановленная таблица становится основой дальнейшего анализа, поэтому важно понимать не только итоговую ошибку метода, но и его устойчивость к шуму, количеству пропусков, структуре наблюдений и выбору ранга.

В экономических данных такая задача выглядит особенно естественно. Показатели стран, регионов, компаний или временных периодов часто зависят от небольшого числа скрытых факторов. Это приводит к гипотезе низкого ранга: большая матрица может приближенно описываться через несколько латентных компонент. На этой идее строится задача \textit{matrix completion}, то есть восстановление недостающих элементов по наблюдаемой части матрицы.

Классические подходы к matrix completion часто используют выпуклую релаксацию через ядерную норму. У таких методов есть сильная теория, но в больших задачах они могут быть тяжелыми вычислительно, поскольку требуют сингулярных разложений и других дорогих операций. На этом фоне риманов градиентный спуск интересен как промежуточный путь: он работает напрямую на множестве матриц фиксированного ранга и использует геометрию этого множества при построении шага.

В этой работе мы рассматриваем тему не только как реализацию одного алгоритма. Основной интерес состоит в том, чтобы понять поведение риманова градиентного спуска в разных режимах: при шуме, разной доле наблюдаемых элементов, разных инициализациях и неверном выборе ранга. Такой подход позволяет сделать из курсовой не просто обзор метода, а небольшой вычислительный исследовательский проект.

\textbf{Цель исследования.} Реализовать и протестировать метод риманова градиентного спуска на многообразии матриц фиксированного ранга для задачи заполнения матриц, а затем сравнить его качество, скорость и устойчивость с базовыми методами matrix completion.

\textbf{Задачи исследования.} Для достижения цели в работе решаются следующие задачи:
\begin{enumerate}[label=\arabic*.]
    \item провести обзор литературы по выпуклым, факторизационным и геометрическим методам заполнения матриц;
    \item описать геометрическую постановку задачи на многообразии матриц фиксированного ранга;
    \item вывести формулу риманова градиента и шаг риманова градиентного спуска для функции невязки по наблюдаемым элементам;
    \item рассмотреть регуляризованную версию риманова спуска для шумных данных;
    \item реализовать экспериментальный код для риманова градиентного спуска, ALS и Soft-Impute как базовых методов сравнения;
    \item подготовить серию синтетических экспериментов с контролируемым рангом, шумом и долей пропусков;
    \item исследовать устойчивость методов к шуму, доле наблюдаемых элементов, выбору ранга и инициализации;
    \item проанализировать скорость сходимости и вычислительную эффективность алгоритмов;
    \item на следующем этапе перейти к компактной реализации через $X=USV^\top$ и прикладной проверке на экономических данных.
\end{enumerate}

\textbf{Объект и предмет исследования.}
\begin{itemize}
    \item \textbf{Объект исследования:} процессы восстановления пропущенных значений в низкоранговых матрицах.
    \item \textbf{Предмет исследования:} риманов градиентный спуск на многообразии матриц фиксированного ранга и его свойства в сравнении с альтернативными алгоритмами matrix completion.
\end{itemize}

\textbf{Методы исследования.} В работе используются методы линейной алгебры, элементы дифференциальной геометрии и римановой оптимизации, анализ литературы, а также вычислительный эксперимент. Качество восстановления оценивается по RMSE, MAE и относительной ошибке в норме Фробениуса, а вычислительная эффективность --- по времени работы и числу итераций.

\textbf{Структура работы.} В первом разделе представлен обзор литературы по matrix completion. Во втором разделе вводятся основные обозначения. В третьем разделе приводится теоретическая часть: сначала кратко вводятся необходимые геометрические понятия, затем формулируется задача на многообразии матриц фиксированного ранга и описывается риманов градиентный спуск. В четвертом разделе описываются реализация и дизайн вычислительных экспериментов. В заключении формулируются промежуточные выводы и дальнейшие шаги работы.

\newpage

\section{Обзор литературы}
\label{sec:literature}

\subsection{Классическая выпуклая линия}

Современная теория matrix completion начинается с выпуклой постановки через ядерную норму. В работе Candes и Recht \cite{candes-recht} показано, что матрицу малого ранга можно восстанавливать по случайно наблюдаемым элементам, если заменить минимизацию ранга на выпуклую релаксацию. Для литературы это был принципиальный шаг: задача перестала быть чисто эвристической и получила строгую математическую основу.

Дальше эта линия была существенно усилена. В работе Candes и Tao \cite{candes-tao} получены почти оптимальные по порядку оценки на число наблюдений, достаточное для восстановления. Статья Recht \cite{recht} дала более компактное и прозрачное доказательство близких результатов. Если смотреть на эти три работы вместе, то они задают базовую выпуклую логику всей темы: сначала показывается, что восстановление вообще возможно, затем уточняются условия, при которых это можно гарантировать.

Для нашей работы эта линия важна по двум причинам. Во-первых, она дает естественный baseline: если мы хотим оценивать новый или более сложный метод, то сравнивать его разумно именно с выпуклой постановкой, у которой есть сильная теория. Во-вторых, эти статьи показывают ограничения такого подхода. Основной акцент в них делается на гарантии восстановления, тогда как вычислительная цена решения в больших задачах остается отдельным вопросом.

\subsection{Алгоритмы для больших матриц}

После теоретических результатов почти сразу возникает вычислительная проблема. Даже если выпуклая модель выглядит естественно, ее надо уметь считать на больших матрицах. В работе Cai, Candes и Shen \cite{cai-svt} предложен алгоритм Singular Value Thresholding, который переводит общую выпуклую постановку в последовательность более управляемых шагов.

Эту же линию продолжает работа Mazumder, Hastie и Tibshirani \cite{mazumder}, где предлагается Soft-Impute. Этот алгоритм особенно важен для прикладной части темы, потому что ориентирован на большие и разреженные матрицы. В отличие от более абстрактного теоретического языка первых статей, здесь уже явно виден переход к реальным вычислительным задачам.

Для нашей работы эти алгоритмы нужны не просто как исторические примеры. Скорее, это естественные конкуренты геометрического подхода. Если риманов градиентный спуск должен быть содержательно интересен, он должен сравниваться не с абстрактной выпуклой моделью, а с реально используемыми численными методами.

\subsection{Шум и более реалистичная постановка}

Классические результаты о точном восстановлении важны, но в прикладной задаче этого недостаточно. В реальных данных шум почти неизбежен: наблюдаемые значения сами по себе измеряются с ошибкой, а шаблон пропусков часто далек от идеального случайного. Работа Candes и Plan \cite{candes-plan} важна именно тем, что переводит matrix completion в постановку с шумом.

Для нашего проекта это ключевой момент. Если ограничиться только безошибочной моделью, то значительная часть будущего сравнения алгоритмов просто теряет смысл. Нас интересует как раз то, что происходит при реальных возмущениях: устойчивость к шуму, качество восстановления скрытых наблюдений и чувствительность метода к меняющимся условиям эксперимента.

\subsection{Невыпуклые и фиксированно-ранговые методы}

Хотя выпуклая линия очень сильна теоретически, она не всегда оптимальна вычислительно. Это естественным образом приводит к невыпуклым постановкам, где ранг задается заранее, а матрица ищется в факторизованной форме. Одной из ранних и важных работ в этом направлении является статья Keshavan, Montanari и Oh \cite{keshavan}, где предложен метод OptSpace.

Для нас OptSpace важен не только как конкретный алгоритм, но и как переходный этап в развитии темы. Здесь исследователь уже не заменяет ранг выпуклой оболочкой, а работает непосредственно с низкоранговой структурой. Именно такая логика и подводит к геометрическим методам на многообразиях фиксированного ранга.

\subsection{Геометрический подход}

Прямой геометрический подход к matrix completion подробно развивается в работе Vandereycken \cite{vandereycken}. В этой статье задача формулируется как оптимизация функции невязки на многообразии матриц фиксированного ранга, а все основные объекты --- касательное пространство, риманов градиент, ретракция --- вводятся в естественной для многообразия форме. Для нашей темы эта статья является одной из центральных, потому что именно она задает базовую схему риманова градиентного метода, который мы реализуем.

Работа Mishra и Sepulchre \cite{r3mc} развивает ту же линию, но уже с более сложной геометрической конструкцией. Авторы предлагают алгоритм R3MC и подчеркивают, что выбор подходящей римановой метрики может заметно повлиять на устойчивость и численное поведение метода. Для нашего проекта это особенно полезно, потому что мы заранее хотим рассматривать не только итоговую ошибку восстановления, но и более тонкие характеристики: зависимость от старта, чувствительность к шуму и скорость сходимости.

Важное отличие геометрической линии от выпуклой состоит в том, что здесь акцент смещается с вопроса ``когда восстановление гарантировано'' на вопрос ``как именно ведет себя алгоритм при прямой оптимизации на низкоранговом множестве''. Именно поэтому геометрические методы хорошо подходят для проекта с исследовательским уклоном.

\subsection{Расширения и прикладной контекст}

Литература по matrix completion давно вышла за рамки одной базовой постановки. В работе Jain и Dhillon \cite{inductive} рассматривается \textit{inductive matrix completion}, где кроме самой матрицы используются и дополнительные признаки строк и столбцов. Это важно, потому что в прикладных данных часто есть внешняя информация, которую можно использовать при восстановлении.

Статья Athey et al. \cite{causal-panel} показывает, что методы matrix completion уже применяются и в задачах с панельными экономическими данными. Для нашей работы это служит прикладным ориентиром: тема не ограничивается абстрактной линейной алгеброй, а реально используется в задачах, близких к экономике и анализу панелей.

Наконец, обзор Davenport и Romberg \cite{overview} полезен тем, что он связывает разные направления low-rank recovery в одну картину. Из него хорошо видно, что выбор метода почти всегда связан с компромиссом между теоретической строгостью, вычислительной стоимостью и гибкостью модели.

\subsection{Выводы из обзора литературы}

Из обзора литературы для нашей работы следуют несколько важных выводов. Во-первых, выпуклые методы с ядерной нормой остаются естественным baseline, потому что они хорошо изучены и понятны теоретически. Во-вторых, методы фиксированного ранга и геометрические алгоритмы особенно интересны в больших задачах, где прямая работа с низкоранговой структурой может дать вычислительное преимущество. В-третьих, именно для геометрических методов особенно важны практические вопросы, которые не сводятся к одной только гарантии восстановления: устойчивость к шуму, зависимость от инициализации, чувствительность к выбору ранга и поведение по ходу итераций.

Именно здесь и находится место нашей работы. Мы не пытаемся доказывать новые теоремы о восстановлении, но и не ограничиваемся простой реализацией известного алгоритма. Основная идея проекта состоит в том, чтобы реализовать риманов градиентный спуск на многообразии матриц фиксированного ранга и систематически сравнить его с альтернативами по нескольким критериям.

\newpage

\section{Основные определения и обозначения}
\label{sec:definitions}

Пусть дана матрица
\[
Y \in \R^{m \times n},
\]
в которой наблюдаются только элементы с индексами из множества
\[
\Omega \subseteq \{1,\dots,m\} \times \{1,\dots,n\}.
\]
Предполагается, что существует скрытая матрица
\[
X_* \in \R^{m \times n},
\]
имеющая малый ранг, и наблюдения удовлетворяют соотношению
\[
Y_{ij} = (X_*)_{ij} + \varepsilon_{ij}, \qquad (i,j)\in\Omega,
\]
где $\varepsilon_{ij}$ описывает шум.

\begin{definition}
Введем маску наблюдений $W\in\R^{m\times n}$:
\[
W_{ij} =
\begin{cases}
1, & (i,j)\in\Omega,\\
0, & (i,j)\notin\Omega.
\end{cases}
\]
Тогда оператор наблюдаемых элементов удобно записывать через адамарово произведение:
\[
\Proj_\Omega(X)=W\odot X.
\]
Здесь $\odot$ означает покомпонентное произведение матриц.
\end{definition}

Тогда задача matrix completion состоит в восстановлении матрицы $X_*$ по наблюдаемой части $\Proj_\Omega(Y)$ при предположении, что искомая матрица имеет малый ранг.

В дальнейшем используются следующие обозначения:
\begin{itemize}
    \item $\fro{X}$ --- норма Фробениуса матрицы $X$;
    \item $\nuc{X}$ --- ядерная норма матрицы $X$;
    \item $\rank(X)$ --- ранг матрицы $X$;
    \item $\Mcal_r = \{X \in \R^{m\times n} : \rank(X)=r\}$ --- множество матриц фиксированного ранга $r$;
    \item $U\Sigma V^\top$ --- компактное сингулярное разложение матрицы ранга $r$.
\end{itemize}

\newpage

\section{Теоретическая часть}
\label{sec:theory}

\subsection{Карты, атласы и гладкость}

В этом разделе используется только тот фрагмент дифференциальной геометрии, который нужен для постановки риманова градиентного спуска. Материал опирается на лекционные и семинарские курсы по дифференциальной геометрии А. О. Иванова, А. Т. Фоменко и А. В. Пенского \cite{ivanov-dg,fomenko-dg,penskoy-seminars}.

\begin{definition}
Пусть $M$ --- топологическое пространство. Картой размерности $d$ называется пара $(U,\varphi)$, где $U\subset M$ --- открытое множество, а
\[
\varphi:U\to \varphi(U)\subset\R^d
\]
--- гомеоморфизм на открытое множество $\varphi(U)$.
\end{definition}

Карта позволяет заменить точки из $U$ их координатами в $\R^d$. Если таких карт достаточно, чтобы покрыть все множество, мы получаем атлас.

\begin{definition}
Набор карт $\mathcal A=\{(U_\alpha,\varphi_\alpha)\}$ называется атласом, если
\[
M=\bigcup_\alpha U_\alpha.
\]
Если все карты имеют значения в $\R^d$, то число $d$ называется размерностью многообразия.
\end{definition}

\begin{definition}
Топологическое пространство $M$ называется $d$-мерным топологическим многообразием, если у каждой точки $p\in M$ есть окрестность, гомеоморфная открытому множеству в $\R^d$.
\end{definition}

Иначе говоря, многообразие может быть устроено глобально сложно, но локально оно выглядит как обычное евклидово пространство. Именно поэтому на нем можно вводить локальные координаты.

Если две карты пересекаются, то одна и та же точка получает две записи. Поэтому важно отображение перехода между координатами:
\[
\psi\circ\varphi^{-1}:\varphi(U\cap V)\to\psi(U\cap V).
\]

\begin{definition}
Атлас называется гладким, если все отображения перехода между его картами гладкие. Множество с таким атласом называется гладким многообразием.
\end{definition}

\begin{definition}
Пусть $M$ и $N$ --- гладкие многообразия. Отображение $F:M\to N$ называется гладким, если в любых картах его координатная запись
\[
\psi\circ F\circ\varphi^{-1}
\]
является гладким отображением между открытыми подмножествами евклидовых пространств.
\end{definition}

Гладкость нам нужна для того чтобы производные, касательные векторы и градиенты не зависели от произвольного выбора координат.

\subsection{Вложенные многообразия и регулярные уровни}

В оптимизации чаще всего встречаются не абстрактные многообразия, а подмножества евклидова пространства, заданные ограничениями. Такие множества называют вложенными многообразиями.

\begin{definition}
Подмножество $M\subset\R^N$ называется вложенным гладким многообразием размерности $d$, если в окрестности каждой своей точки оно после гладкой замены координат выглядит как $\R^d\times\{0\}^{N-d}$.
\end{definition}

Один из стандартных способов получить вложенное многообразие --- задать его системой уравнений. Для этого нужны регулярные точки.

\begin{definition}
Пусть $F:\R^N\to\R^k$ --- гладкое отображение. Точка $x$ называется регулярной точкой, если матрица Якоби $DF(x)$ имеет полный ранг $k$. Если ранг меньше $k$, точка называется особой.
\end{definition}

\begin{definition}
Точка $y\in\R^k$ называется регулярным значением отображения $F$, если каждая точка из прообраза $F^{-1}(y)$ является регулярной. Если в прообразе есть хотя бы одна особая точка, то $y$ называется особым значением.
\end{definition}

\begin{theorem}[Теорема о регулярном уровне]
Пусть $F:\R^N\to\R^k$ --- гладкое отображение, и для всех $x$ таких, что $F(x)=0$, матрица $DF(x)$ имеет ранг $k$. Тогда множество
\[
M=\{x\in\R^N:F(x)=0\}
\]
является вложенным гладким многообразием размерности $N-k$.
\end{theorem}

Это утверждение является прямым следствием теоремы о неявной функции: если уравнения локально независимы, то часть координат можно выразить через остальные. Поэтому остается $N-k$ свободных координат.

Обратная идея тоже важна: локально вложенное многообразие можно описывать системой гладких уравнений с полным рангом Якоби. Для нашей темы это объясняет, почему ограничения в оптимизации можно изучать геометрически, а не только как набор формальных равенств.

\subsection{Касательные векторы}

Касательный вектор сначала можно определить через координаты. Пусть $x\in M$ и выбрана карта $(U,\varphi)$ около $x$. Тогда вектор в координатах --- это элемент $v\in\R^d$. Если перейти к другой карте $(V,\psi)$, тот же вектор должен измениться по правилу
\[
v_\psi =
D(\psi\circ\varphi^{-1})(\varphi(x))\,v_\varphi.
\]
Поэтому определение касательного вектора не зависит от выбора конкретной карты: при смене координат его координатное представление меняется по правилу дифференциала отображения перехода.

Для наших целей удобнее эквивалентное определение через кривые:
\[
T_xM=\{\dot\gamma(0):\gamma(t)\in M,\ \gamma(0)=x\}.
\]
Существует стандартная теорема об эквивалентности координатного определения и определения через кривые; ее доказательство здесь не приводится. В дальнейшем мы будем пользоваться определением через кривые, потому что оно удобнее для вложенных многообразий.

\subsection{Риманова метрика и градиент}

\begin{definition}
Риманова метрика --- это выбор скалярного произведения $\langle\cdot,\cdot\rangle_x$ на каждом касательном пространстве $T_xM$, гладко зависящий от точки $x$. Для любых $\xi,\eta,\zeta\in T_xM$ и чисел $a,b$ оно удовлетворяет:
\[
\langle a\xi+b\eta,\zeta\rangle_x
=a\langle\xi,\zeta\rangle_x+b\langle\eta,\zeta\rangle_x,
\]
\[
\langle \xi,\eta\rangle_x=\langle\eta,\xi\rangle_x,\qquad
\langle\xi,\xi\rangle_x>0\quad \text{при }\xi\ne 0.
\]
\end{definition}

Метрика нужна для того, чтобы измерять длины и углы касательных векторов. Без нее нельзя однозначно сказать, какое касательное направление является направлением наибольшего возрастания функции.

В этой работе используется метрика, индуцированная нормой Фробениуса. Для матриц одинакового размера:
\[
\langle A,B\rangle=\tr(A^\top B).
\]
Соответствующая норма:
\[
\fro{A}=\sqrt{\langle A,A\rangle}.
\]
Такой выбор естественен, потому что сама функция потерь в задаче matrix completion записывается через норму Фробениуса.

Пусть $f:M\to\R$ --- гладкая функция. Ее дифференциал $Df(x)$ является линейной функцией на $T_xM$: он принимает касательное направление $\xi$ и возвращает производную функции вдоль этого направления.

\begin{definition}
Риманов градиент $\grad f(x)$ определяется условием
\[
Df(x)[\xi]=\langle \grad f(x),\xi\rangle_x
\qquad \text{для всех }\xi\in T_xM.
\]
\end{definition}

Таким образом, градиент определяется не только функцией, но и выбранной метрикой: именно метрика отождествляет линейную форму $Df(x)$ с касательным вектором.

Если $M$ вложено в евклидово пространство и метрика индуцирована обычным скалярным произведением, то риманов градиент получается проекцией евклидова градиента на касательное пространство:
\[
\grad_M f(x)=\Proj_{T_xM}(\nabla f(x)).
\]
Нормальная компонента евклидова градиента не задает допустимое движение вдоль многообразия. Поэтому в алгоритме остается только касательная часть. Именно эта формула будет использоваться дальше для матриц фиксированного ранга.

\subsection{Геодезические и ретракции}

На многообразии градиент $\grad f(x)$ лежит в касательном пространстве $T_xM$, а не в самом многообразии. Поэтому выражение
\[
x-\eta \grad f(x)
\]
обычно не является точкой из $M$. Нужно сначала сделать шаг в касательном направлении, а затем вернуться на допустимое множество.

Самый геометрически точный способ --- двигаться по геодезической. Геодезическая играет роль прямой линии на многообразии: ее начальная точка равна $x$, а начальная скорость равна выбранному касательному вектору. Если обозначить через $\operatorname{Exp}_x$ экспоненциальное отображение, то идеальный шаг можно записать так:
\[
x_{k+1}=\operatorname{Exp}_{x_k}\bigl(-\eta_k\grad f(x_k)\bigr).
\]
Проблема в том, что экспоненциальное отображение часто трудно вычислять явно. Поэтому в алгоритмах обычно используют более простую замену --- ретракцию.

\begin{definition}
Ретракцией называется отображение $\Retr_x:T_xM\to M$, которое удовлетворяет двум условиям:
\[
\Retr_x(0_x)=x,
\qquad
D\Retr_x(0_x)=\operatorname{id}_{T_xM}.
\]
\end{definition}

Первое условие говорит, что нулевой шаг оставляет точку на месте. Второе означает, что при малых шагах ретракция не искажает касательное направление в первом порядке. Поэтому ретракция может заменить геодезическую в градиентном методе: она сохраняет нужную локальную геометрию, но обычно считается проще.

Для матриц фиксированного ранга естественная ретракция будет устроена так: сделать касательный шаг, а затем заменить полученную матрицу ближайшей матрицей ранга $r$. Ниже это приведет к усеченному SVD.

\subsection{Постановка задачи}

В задаче заполнения матриц наблюдается только часть элементов матрицы $Y$. Классическая выпуклая постановка использует ядерную норму:
\begin{equation}
\label{eq:convex-final}
\min_{X \in \R^{m\times n}}
\frac12 \fro{\Proj_\Omega(X-Y)}^2 + \lambda \nuc{X}.
\end{equation}
Первое слагаемое отвечает за согласование с наблюдаемыми элементами, второе поощряет малый ранг. Эта модель выпуклая, но в больших задачах может быть вычислительно дорогой.

В fixed-rank подходе ранг задается заранее, и оптимизация ведется напрямую по множеству матриц ранга $r$:
\begin{equation}
\label{eq:fixed-rank-final}
\min_{X \in \Mcal_r}
\frac12 \fro{\Proj_\Omega(X-Y)}^2,
\qquad
\Mcal_r=\{X\in\R^{m\times n}:\rank(X)=r\}.
\end{equation}

Для шумных данных дополнительно рассматривается L2-регуляризация:
\begin{equation}
\label{eq:fixed-rank-regularized}
\min_{X \in \Mcal_r} f_\lambda(X)
= \frac12 \fro{\Proj_\Omega(X-Y)}^2 + \frac{\lambda}{2}\fro{X}^2,
\qquad \lambda\ge 0.
\end{equation}
Штраф $\fro{X}^2$ ограничивает размер решения и снижает риск подстройки под шум.

\subsection{Многообразие матриц фиксированного ранга}

Покажем, что область оптимизации в \eqref{eq:fixed-rank-final} является гладким многообразием.

\begin{proposition}
Множество
\[
\Mcal_r=\{X\in\R^{m\times n}:\rank(X)=r\}
\]
является гладким вложенным многообразием размерности
\[
\dim \Mcal_r=r(m+n-r).
\]
\end{proposition}

\begin{proof}
Возьмем матрицу $X\in\Mcal_r$. У нее есть невырожденный минор размера $r\times r$. После перестановки строк и столбцов можно считать, что он находится в левом верхнем углу. Тогда в окрестности $X$ матрицу можно записать блоками
\[
X=
\begin{pmatrix}
A & B\\
C & D
\end{pmatrix},
\qquad A\in\R^{r\times r},\quad \det A\ne 0.
\]
Пока блок $A$ невырожден, условие $\rank(X)=r$ локально эквивалентно равенству
\[
D=CA^{-1}B.
\]
Следовательно, свободными параметрами являются блоки $A$, $B$ и $C$, а блок $D$ гладко выражается через них. Число свободных параметров равно
\[
r^2+r(n-r)+(m-r)r=r(m+n-r).
\]
Эти параметры задают локальные координаты на $\Mcal_r$. При переходе к другой ненулевой минорной карте координаты меняются гладко, поэтому $\Mcal_r$ является гладким вложенным многообразием.
\end{proof}

Множество $\{X:\rank(X)\le r\}$ уже не является гладким многообразием всюду: в точках меньшего ранга возникают особенности. Поэтому дальше рассматривается именно страт $\Mcal_r=\{X:\rank(X)=r\}$.

Пусть
\[
X=U\Sigma V^\top
\]
--- компактное SVD, где $U^\top U=I_r$, $V^\top V=I_r$, а $\Sigma$ невырождена. Тогда касательное пространство можно записать в виде
\begin{equation}
\label{eq:tangent-space-final}
T_X\Mcal_r =
\left\{
UMV^\top + U_pV^\top + UV_p^\top :
M \in \R^{r\times r},\,
U^\top U_p = 0,\,
V^\top V_p = 0
\right\}.
\end{equation}
Первое слагаемое соответствует изменению центрального блока, второе --- столбцового подпространства, третье --- строкового подпространства.

Для вычисления риманова градиента нужна проекция на касательное пространство. Если $P_U=UU^\top$ и $P_V=VV^\top$, то
\begin{equation}
\label{eq:tangent-projection-final}
\Proj_{T_X\Mcal_r}(Z)=P_UZ+ZP_V-P_UZP_V.
\end{equation}
Эта формула получается из разложения $Z$ на четыре части:
\[
Z=P_UZP_V+P_UZ(I-P_V)+(I-P_U)ZP_V+(I-P_U)Z(I-P_V).
\]
Последнее слагаемое является нормальной компонентой; после его удаления получается \eqref{eq:tangent-projection-final}.

\subsection{Риманов градиентный спуск}

Рассмотрим функцию
\[
f(X)=\frac12\fro{\Proj_\Omega(X-Y)}^2.
\]
Ее евклидов градиент равен
\[
\nabla f(X)=\Proj_\Omega(X-Y).
\]
Для регуляризованной функции \eqref{eq:fixed-rank-regularized} получаем
\[
\nabla f_\lambda(X)=\Proj_\Omega(X-Y)+\lambda X.
\]
Это следует из дифференцирования по произвольному приращению $H$:
\[
\left.\frac{d}{dt}\right|_{t=0}
\frac12\fro{\Proj_\Omega(X+tH-Y)}^2
=
\langle \Proj_\Omega(X-Y),\Proj_\Omega(H)\rangle.
\]
Так как $\Proj_\Omega$ является ортогональной проекцией,
\[
\langle \Proj_\Omega(X-Y),\Proj_\Omega(H)\rangle
=
\langle \Proj_\Omega(X-Y),H\rangle.
\]
Отсюда получается формула для $\nabla f(X)$, а регуляризация добавляет градиент $\lambda X$.

Риманов градиент --- это проекция евклидова градиента на касательное пространство:
\begin{equation}
\label{eq:riemannian-gradient-final}
\grad_{\Mcal_r} f_\lambda(X)
=
\Proj_{T_X\Mcal_r}\bigl(\Proj_\Omega(X-Y)+\lambda X\bigr).
\end{equation}

После касательного шага нужно вернуться на $\Mcal_r$. Для этого используется ретракция:
\begin{equation}
\label{eq:fixed-rank-retraction-final}
\Retr_X(\Xi)=\Pi_r(X+\Xi),
\end{equation}
где $\Pi_r$ --- усеченное SVD. Если
\[
X+\Xi=\sum_i \sigma_i u_i v_i^\top
\]
--- сингулярное разложение, то
\[
\Pi_r(X+\Xi)=\sum_{i=1}^r \sigma_i u_i v_i^\top.
\]
По теореме Эккарта--Янга это ближайшая матрица ранга не выше $r$ в норме Фробениуса. В окрестности точки ранга $r$ такая операция возвращает нас на $\Mcal_r$.

Проверим условия ретракции. При $\Xi=0$ имеем $\Pi_r(X)=X$. Кроме того, для малого $\Xi\in T_X\Mcal_r$:
\[
\Pi_r(X+\Xi)=X+\Xi+O(\fro{\Xi}^2).
\]
Следовательно, дифференциал ретракции в нуле является тождественным на $T_X\Mcal_r$. Геодезики здесь не используются: для градиентного метода достаточно совпадения с касательным шагом в первом порядке, а усеченное SVD считается проще.

Итоговый шаг риманова градиентного спуска:
\begin{equation}
\label{eq:rgd-final}
X_{k+1}
=
\Retr_{X_k}\bigl(-\eta_k\grad_{\Mcal_r} f_\lambda(X_k)\bigr).
\end{equation}
Длина шага $\eta_k$ подбирается фиксированно или с помощью backtracking line search.

\begin{algorithm}[H]
\caption{Риманов градиентный спуск на $\Mcal_r$}
\label{alg:rgd-final}
\begin{algorithmic}[1]
\Require наблюдаемая матрица $Y$, маска $\Omega$, ранг $r$, параметр $\lambda$, начальная матрица $X_0\in\Mcal_r$
\Ensure приближение $\widehat X$ ранга $r$
\For{$k=0,1,2,\ldots$}
    \State $G_k \gets \Proj_\Omega(X_k-Y)+\lambda X_k$
    \State $\xi_k \gets \Proj_{T_{X_k}\Mcal_r}(G_k)$
    \State выбрать шаг $\eta_k$
    \State $X_{k+1}\gets \Pi_r(X_k-\eta_k\xi_k)$
    \If{критерий остановки выполнен}
        \State \textbf{break}
    \EndIf
\EndFor
\State \Return $X_k$
\end{algorithmic}
\end{algorithm}

\subsection{Компактная реализация через факторы}

На текущем этапе основной алгоритм записан через матрицу $X$ и ретракцию с помощью усеченного SVD. Это удобно для понимания и для первой реализации, но не идеально для больших матриц: на каждой итерации приходится работать с матрицей размера $m\times n$.

Следующий естественный шаг --- хранить текущую точку в виде
\[
X=USV^\top,
\qquad
U^\top U=I_r,\quad V^\top V=I_r.
\]
Идея состоит в том, чтобы выполнять риманов шаг не как полный переход к большой матрице и обратно, а через факторы $U$, $S$, $V$ и малые матрицы размера порядка $r$. Такой вариант должен быть вычислительно выгоднее, особенно когда $r\ll m,n$.

При этом важно, что это не просто обычный градиентный спуск по независимым переменным $U$, $S$, $V$. Нужно сохранить геометрию многообразия фиксированного ранга: касательное направление, проекцию и возврат на множество матриц ранга $r$. Полная схема такого шага сейчас находится в процессе уточнения. Предполагается, что регуляризация из \eqref{eq:fixed-rank-regularized} переносится в этот вариант так же: в евклидовом градиенте появляется добавка $\lambda X$.

В финальной версии работы компактная реализация будет рассмотрена как основной путь к масштабированию метода. В предварительном тексте мы фиксируем ее идею, но основные эксперименты ниже пока интерпретируем через более прозрачный вариант с SVD-ретракцией.

\newpage

\section{Практическая часть}
\label{sec:empirical}

\subsection{Код и постановка эксперимента}

Для практической части написан отдельный Python-скрипт:
\begin{center}
\path{matrix_completion_experiments.py}
\end{center}
Скрипт генерирует низкоранговые матрицы, скрывает часть элементов, добавляет шум и сравнивает методы восстановления. Такой синтетический эксперимент нужен как контролируемая проверка: мы знаем настоящую матрицу $X_*$ и можем честно измерять ошибку.

Матрица строится по модели
\[
X_* = AB^\top,
\]
где $A\in\R^{m\times r}$ и $B\in\R^{n\times r}$. Затем часть элементов отдается алгоритмам для обучения, а часть известных элементов откладывается как test. На наблюдаемых позициях берется
\[
Y_{ij}=(X_*)_{ij}+\varepsilon_{ij}, \qquad (i,j)\in\Omega,
\]
где $\varepsilon_{ij}$ отвечает за шум.

\subsection{Методы и метрики}

В текущей версии сравниваются:
\begin{enumerate}[label=\arabic*.]
    \item ALS как простой факторизационный baseline;
    \item Soft-Impute как baseline, связанный с ядерной нормой;
    \item риманов градиентный спуск без регуляризации;
    \item риманов градиентный спуск с L2-регуляризацией;
    \item риманов градиентный спуск с выбором $\lambda$ по validation-множеству.
\end{enumerate}

Основные метрики:
\begin{enumerate}[label=\arabic*.]
    \item RMSE и MAE на test-элементах;
    \item относительная ошибка $\fro{\widehat X-X_*}/\fro{X_*}$;
    \item время работы;
    \item число итераций.
\end{enumerate}

Параметр $\lambda$ подбирается двумя способами. Сначала мы смотрим сетку значений, чтобы понять, как регуляризация влияет на качество. Затем используем честный validation-подбор: часть наблюдаемых элементов временно откладывается, $\lambda$ выбирается по ним, а итоговая модель заново обучается на всех наблюдаемых элементах.

\begin{algorithm}[H]
\caption{Validation-подбор параметра $\lambda$}
\label{alg:validation-lambda-final}
\begin{algorithmic}[1]
\Require наблюдаемые индексы $\Omega_{\mathrm{train}}$, сетка $\Lambda$, доля validation $\rho$
\Ensure выбранное значение $\lambda_*$ и итоговая матрица $\widehat X$
\State разбить $\Omega_{\mathrm{train}}$ на $\Omega_{\mathrm{fit}}$ и $\Omega_{\mathrm{val}}$
\For{$\lambda\in\Lambda$}
    \State обучить RGD на $\Omega_{\mathrm{fit}}$ с параметром $\lambda$
    \State посчитать $\operatorname{RMSE}_{\mathrm{val}}(\lambda)$ на $\Omega_{\mathrm{val}}$
\EndFor
\State $\lambda_*\gets \argmin_{\lambda\in\Lambda}\operatorname{RMSE}_{\mathrm{val}}(\lambda)$
\State заново обучить RGD с $\lambda_*$ на всех индексах $\Omega_{\mathrm{train}}$
\State оценить итоговую ошибку на test-множестве
\State \Return $\lambda_*$ и $\widehat X$
\end{algorithmic}
\end{algorithm}

\subsection{Предварительные результаты}

В текущем запуске использовались синтетические матрицы размера $60\times 60$ истинного ранга $4$. Каждый сценарий запускался в трех повторениях. Менялись уровень шума, доля наблюдаемых элементов, тип инициализации и ранг, который передается алгоритму.

\begin{table}[H]
\centering
\small
\caption{Средние результаты по всем сценариям}
\label{tab:main-results-final}
\begin{tabular}{lrrrr}
\toprule
Метод & RMSE & Отн. ошибка & Время, c & Итерации \\
\midrule
ALS & 0.0720 & 0.7307 & 0.1281 & 49.1 \\
Soft-Impute & 0.0420 & 0.4191 & 0.0952 & 110.3 \\
RGD & 0.0443 & 0.4393 & 0.1708 & 195.0 \\
RGD, $\lambda=10^{-2}$ & 0.0386 & 0.3820 & 0.1733 & 190.1 \\
RGD, $\lambda=3\cdot10^{-2}$ & 0.0390 & 0.3850 & 0.1520 & 182.8 \\
RGD, validation $\lambda$ & 0.0359 & 0.3571 & 1.3552 & 1600.4 \\
\bottomrule
\end{tabular}
\end{table}

По таблице \ref{tab:main-results-final} видно, что обычный RGD близок к Soft-Impute, но не всегда его обгоняет. После добавления L2-регуляризации качество заметно улучшается. Лучший результат дает validation-подбор $\lambda$, но он намного дороже по времени, потому что для одного сценария приходится обучать несколько моделей.

\begin{table}[H]
\centering
\small
\caption{Средний RMSE при разных уровнях шума}
\label{tab:noise-results-final}
\begin{tabular}{lrrrr}
\toprule
Шум & ALS & Soft-Impute & RGD & RGD, validation $\lambda$ \\
\midrule
0.00 & 0.0450 & 0.0367 & 0.0263 & 0.0240 \\
0.02 & 0.0601 & 0.0395 & 0.0368 & 0.0328 \\
0.05 & 0.1110 & 0.0498 & 0.0698 & 0.0508 \\
\bottomrule
\end{tabular}
\end{table}

Таблица \ref{tab:noise-results-final} показывает главный вывод по шуму. При нулевом и умеренном шуме риманов подход выглядит очень хорошо. При сильном шуме Soft-Impute остается более устойчивым. Это логично: метод с ядерной нормой сильнее сглаживает сингулярные значения, а fixed-rank метод может переобучаться под шум, если регуляризация подобрана недостаточно удачно.

\begin{table}[H]
\centering
\small
\caption{Сетка L2-регуляризации для RGD}
\label{tab:l2-sweep-final}
\begin{tabular}{rrrr}
\toprule
$\lambda$ & RMSE & Отн. ошибка & Время, c \\
\midrule
$10^{-4}$ & 0.0442 & 0.4379 & 0.1730 \\
$3\cdot10^{-4}$ & 0.0439 & 0.4352 & 0.1783 \\
$10^{-3}$ & 0.0430 & 0.4267 & 0.1802 \\
$3\cdot10^{-3}$ & 0.0412 & 0.4085 & 0.1807 \\
$10^{-2}$ & 0.0386 & 0.3820 & 0.1733 \\
$3\cdot10^{-2}$ & 0.0390 & 0.3850 & 0.1520 \\
$10^{-1}$ & 0.0463 & 0.4564 & 0.1434 \\
\bottomrule
\end{tabular}
\end{table}

Из таблицы \ref{tab:l2-sweep-final} видно, что слишком слабая регуляризация почти не меняет метод, а слишком сильная уже портит качество. В текущем эксперименте лучшие значения находятся около $10^{-2}$ и $3\cdot10^{-2}$.

\begin{table}[H]
\centering
\small
\caption{Как часто validation выбирал разные значения $\lambda$}
\label{tab:selected-lambda-final}
\begin{tabular}{rrr}
\toprule
$\lambda$ & Число выборов & Доля \\
\midrule
$10^{-4}$ & 15 & 9.3\% \\
$3\cdot10^{-4}$ & 4 & 2.5\% \\
$10^{-3}$ & 2 & 1.2\% \\
$3\cdot10^{-3}$ & 30 & 18.5\% \\
$10^{-2}$ & 44 & 27.2\% \\
$3\cdot10^{-2}$ & 42 & 25.9\% \\
$10^{-1}$ & 25 & 15.4\% \\
\bottomrule
\end{tabular}
\end{table}

Validation чаще всего выбирает значения порядка $10^{-2}$ и $3\cdot10^{-2}$, но иногда выбирает и более сильную регуляризацию. Поэтому нельзя просто один раз зафиксировать $\lambda$ и считать, что задача решена: сила регуляризации зависит от шума, доли наблюдений и выбранного ранга.

\subsection{Что это дает для дальнейшей работы}

Предварительные эксперименты поддерживают следующую рабочую гипотезу: риманов градиентный спуск полезен, когда данные действительно близки к низкоранговым и шум не слишком большой. L2-регуляризация заметно улучшает устойчивость, особенно если ранг выбран не идеально.

При этом текущая SVD-ретракция остается вычислительно дорогой для больших матриц. Поэтому следующий шаг --- компактная реализация через $X=USV^\top$. Ожидается, что она будет более оптимальной по времени, потому что вместо полного SVD большой матрицы можно будет работать с факторами и малыми матрицами. Регуляризация при этом сохранится: в градиенте по-прежнему будет добавка $\lambda X$.

После синтетических экспериментов нужно перейти к экономическим или панельным данным. Часть известных значений будет скрываться искусственно, чтобы методы можно было сравнить по восстановлению этих значений. При этом реальные пропуски могут быть неслучайными, поэтому прикладной эксперимент лучше рассматривать как дополнительную проверку, а не как замену синтетике.

\newpage

\section*{Заключение}
\phantomsection
\addcontentsline{toc}{section}{Заключение}

\textbf{Выводы.} В предварительном тексте сформирована основа курсовой работы по геометрическим методам для заполнения матриц. Во введении уточнена исследовательская постановка, в обзоре литературы показан переход от выпуклых методов к fixed-rank и римановым подходам, а в теоретической части добавлены необходимые элементы дифференциальной геометрии: касательное пространство, риманов градиент, проекция и ретракция.

\textbf{Основной результат текущего этапа.} Работа уже имеет не только теоретический, но и практический каркас. Подготовлен код для синтетических экспериментов, где можно отдельно менять шум, долю наблюдений, инициализацию, выбранный ранг и уровень L2-регуляризации в римановом методе. Первые результаты показывают, что регуляризация заметно улучшает качество RGD, а validation-подбор $\lambda$ позволяет честно выбирать параметр без подглядывания в test-множество.

\textbf{Дальнейшие шаги.} На следующем этапе необходимо уточнить и реализовать компактную версию риманова спуска через $X=USV^\top$, подготовить графики сходимости и перейти к проверке на прикладных экономических данных. После этого можно будет формулировать уже не только методические, но и исследовательские выводы о режимах, в которых риманов градиентный спуск оказывается наиболее полезным.

\begin{thebibliography}{99}
    \phantomsection
    \addcontentsline{toc}{section}{Список литературы}

\bibitem{ivanov-dg}
А. О. Иванов.
\newblock Классическая дифференциальная геометрия; Дифференциальная геометрия и топология.
\newblock Лекционные материалы курса, мехмат МГУ / Teach-in.
\newblock URL: \url{https://teach-in.ru/course/differential-geometry-and-tensor-analysis-ivanov};
\url{https://teach-in.ru/course/differential-geometry-and-tensor-analysis-ivanov-p2}.

\bibitem{fomenko-dg}
А. Т. Фоменко.
\newblock Классическая дифференциальная геометрия.
\newblock Лекционные материалы курса, мехмат МГУ / Teach-in.
\newblock URL: \url{https://teach-in.ru/course/classical-difgeom-fomenko/about}.

\bibitem{penskoy-seminars}
А. В. Пенской.
\newblock Классическая дифференциальная геометрия: семинары.
\newblock Семинарские материалы курса, мехмат МГУ / Teach-in.
\newblock URL: \url{https://teach-in.ru/course/kdg-seminars-penskoy/about}.

\bibitem{candes-recht}
Emmanuel J. Cand\`es, Benjamin Recht.
\newblock Exact Matrix Completion via Convex Optimization.
\newblock \textit{arXiv preprint} arXiv:0805.4471, 2008.
\newblock URL: \url{https://arxiv.org/abs/0805.4471}.

\bibitem{cai-svt}
Jian-Feng Cai, Emmanuel J. Cand\`es, Zuowei Shen.
\newblock A Singular Value Thresholding Algorithm for Matrix Completion.
\newblock \textit{arXiv preprint} arXiv:0810.3286, 2008.
\newblock URL: \url{https://arxiv.org/abs/0810.3286}.

\bibitem{candes-tao}
Emmanuel J. Cand\`es, Terence Tao.
\newblock The Power of Convex Relaxation: Near-Optimal Matrix Completion.
\newblock \textit{arXiv preprint} arXiv:0903.1476, 2009.
\newblock URL: \url{https://arxiv.org/abs/0903.1476}.

\bibitem{recht}
Benjamin Recht.
\newblock A Simpler Approach to Matrix Completion.
\newblock \textit{Journal of Machine Learning Research}, 12:3413--3430, 2011.
\newblock URL: \url{https://www.jmlr.org/papers/v12/recht11a.html}.

\bibitem{keshavan}
Raghunandan H. Keshavan, Andrea Montanari, Sewoong Oh.
\newblock Matrix Completion from a Few Entries.
\newblock \textit{arXiv preprint} arXiv:0901.3150, 2009.
\newblock URL: \url{https://arxiv.org/abs/0901.3150}.

\bibitem{candes-plan}
Emmanuel J. Cand\`es, Yaniv Plan.
\newblock Matrix Completion with Noise.
\newblock \textit{arXiv preprint} arXiv:0903.3131, 2009.
\newblock URL: \url{https://arxiv.org/abs/0903.3131}.

\bibitem{mazumder}
Rahul Mazumder, Trevor Hastie, Robert Tibshirani.
\newblock Spectral Regularization Algorithms for Learning Large Incomplete Matrices.
\newblock \textit{Journal of Machine Learning Research}, 11:2287--2322, 2010.
\newblock URL: \url{https://www.jmlr.org/papers/v11/mazumder10a.html}.

\bibitem{inductive}
Prateek Jain, Inderjit S. Dhillon.
\newblock Provable Inductive Matrix Completion.
\newblock \textit{arXiv preprint} arXiv:1306.0626, 2013.
\newblock URL: \url{https://arxiv.org/abs/1306.0626}.

\bibitem{overview}
Mark A. Davenport, Justin Romberg.
\newblock An Overview of Low-Rank Matrix Recovery from Incomplete Observations.
\newblock \textit{arXiv preprint} arXiv:1601.06422, 2016.
\newblock URL: \url{https://arxiv.org/abs/1601.06422}.

\bibitem{vandereycken}
Bart Vandereycken.
\newblock Low-rank Matrix Completion by Riemannian Optimization --- extended version.
\newblock \textit{arXiv preprint} arXiv:1209.3834, 2012.
\newblock URL: \url{https://arxiv.org/abs/1209.3834}.

\bibitem{r3mc}
B. Mishra, R. Sepulchre.
\newblock R3MC: A Riemannian Three-Factor Algorithm for Low-Rank Matrix Completion.
\newblock \textit{arXiv preprint} arXiv:1306.2672, 2013.
\newblock URL: \url{https://arxiv.org/abs/1306.2672}.

\bibitem{causal-panel}
Susan Athey, Mohsen Bayati, Nikolay Doudchenko, Guido W. Imbens, Khashayar Khosravi.
\newblock Matrix Completion Methods for Causal Panel Data Models.
\newblock \textit{arXiv preprint} arXiv:1710.10251, 2017.
\newblock URL: \url{https://arxiv.org/abs/1710.10251}.

\end{thebibliography}

\end{document}
